\section{Related Work}
\label{sec:relatedwork}

\subsection{Patient-specific ECG adaptation}

Personalizing a heartbeat classifier to an individual is a long-standing idea.
The mixture-of-experts classifier of Hu et al.\ combines a global model with a
small patient-specific expert trained on a brief annotated segment
\cite{hu1997patient}, and Kiranyaz et al.\ train a compact 1-D convolutional
network per patient from a few hundred labelled beats \cite{kiranyaz2016realtime}.
Both establish the central empirical fact we build on: a short labelled prefix of
a patient's own record is worth a great deal, particularly for supraventricular
beats that inter-patient models systematically miss. Neither, however, bounds the
\emph{state} that personalization costs. Their patient models are stored and
trained without a byte budget, and neither considers what must be retained
alongside the parameters --- gradients, optimizer moments, and rehearsal data ---
which is precisely the quantity that decides whether personalization can run on a
microcontroller. Large-scale ECG networks \cite{hannun2019cardiologist} sharpen
the same gap in the opposite direction: accuracy at a parameter scale far beyond
an embedded budget.

\subsection{Inter-patient ECG classification}

The MIT--BIH Arrhythmia Database remains a common benchmark for heartbeat
classification \cite{moody2001impact,goldberger2000physiobank}. Patient-wise
evaluation is substantially more demanding than random beat-level splitting,
because morphology correlations within a record inflate optimistic estimates
when beats from the same patient appear in both train and test. The de Chazal
protocol partitions records into disjoint patient groups (DS1/DS2) and combines
morphology with heartbeat-interval features \cite{dechazal2004automatic}. We
adopt the conventional DS1/DS2 record partition, correct its 201/202 cross-split
subject overlap, and use the AAMI class mapping
\cite{ansi2012aami} but study \emph{supervised adaptation} within an unseen
patient's own chronological record, rather than a single static inter-patient
model. Feature-selection work driven explicitly by cross-database generalization
\cite{llamedo2011heartbeat,mar2011optimization} shows how fragile inter-patient
performance is when the evaluation database changes, which motivates our external
validation on INCART and the MIT--BIH Supraventricular Arrhythmia Database
\cite{greenwald1990svdb}.

\subsection{Latent replay under embedded constraints}

Latent replay stores features extracted by a frozen prefix of a network
\cite{pellegrini2020latent}. Quantized latent replay has been demonstrated on
TinyML platforms and motivates joint selection of replay precision and replay
location \cite{ravaglia2021tinyml}. On-device training within a few hundred
kilobytes has also been shown for vision backbones \cite{lin2022ondevice}. These
results largely concern convolutional architectures whose feature maps shrink
with depth. Our replay-point cliff is architecture-specific: transformer tokens
can be several times \emph{larger} than the raw ECG window, so ``replay deeper
to save memory'' does not hold.

Budget-aware replay with an explicit plasticity trade-off has recently appeared
for vision continual learning \cite{zhang2026amrecl}, and memory--latency--accuracy
trade-offs for latent replay have been characterized on extreme-edge RISC-V nodes
\cite{ravaglia2020memlatacc} from the same lineage as our base model's platform.
Our study differs in three ways: the replay-cost non-monotonicity is a property
specific to an expand-then-pool ECG tokenizer rather than a general depth trend,
the setting is patient-specific ECG adaptation rather than class-incremental
vision, and the constraint is a byte-exact persistent-state budget rather than a
bit-precision or feature-map allocation.

\subsection{Source-free and test-time adaptation}

A parallel line of work adapts a trained model to a shifted target distribution
without labels, either by entropy minimization over test batches
\cite{wang2021tent} or by source-hypothesis transfer \cite{liang2020shot}. These
methods are attractive on a device because they need no annotation, but they
adapt statistics and affine parameters rather than allocating a fixed memory
budget, and they do not address the class-imbalance problem that dominates
minority-class ECG performance. Our setting is deliberately the supervised one:
we assume a short labelled prefix, which is the assumption under which the
patient-specific ECG literature operates, and ask how to spend bytes given it.

\subsection{Continual-learning regularization}

Elastic weight consolidation penalizes changes to parameters estimated to be
important for previous tasks \cite{kirkpatrick2017overcoming}. Under our
accounting, storing a full-model FP32 parameter anchor plus a Fisher value
requires $6643\times8=53{,}144$ bytes, already above the nominal 16\,KiB
ceiling before any replay. The same arithmetic applies to synaptic
intelligence \cite{zenke2017si}; distillation-based approaches such as learning
without forgetting \cite{li2018lwf} avoid the anchor but require a copy of the
source model's outputs, and gradient-projection methods \cite{lopezpaz2017gem}
still keep an episodic memory. This establishes infeasibility only for
\emph{full-model} FP32 regularization under the stated assumptions; restricted or
compressed regularization remains a valid future baseline and we do not claim
that regularization is universally infeasible. Work on very small episodic
memories \cite{chaudhry2019tiny} suggests that the replay side of this trade-off
degrades gracefully, which is consistent with what we measure.

Regularization has also been brought inside the low-rank setting directly.
Zheng et al.\ propose \textbf{EWC-LoRA} \cite{zheng2026ewclora}, which
regularizes a \emph{shared} low-rank update with a Fisher-based importance
estimate, keeping storage and inference cost constant as tasks accumulate. That
is the closest published method to the regularization arm of our study, and the
comparison is worth stating precisely rather than by omission. Zheng et al.\
study Fisher-based regularization of shared low-rank updates in
\emph{task-sequential} vision and language continual learning. Our setting
instead asks how replay and trainable depth should share a \emph{kilobyte-scale
persistent-state budget} for patient-specific ECG adaptation, where the binding
constraint is bytes rather than task count. \textbf{We do not implement
EWC-LoRA}, and therefore restrict our regularization comparison to the tested
$B$-factor anchor (Section~\ref{sec:reg}); we make no claim about how EWC-LoRA
would perform under a 16\,KiB ceiling, and note that a Fisher estimate is itself
a persistent-state cost that would have to be charged against the budget.

\subsection{Parameter-efficient adaptation}

LoRA freezes the original weight matrix and learns a low-rank update
\cite{hu2022lora}. Although introduced for large language models, the mechanism
is equally useful when the goal is to preserve a tiny inference model while
exposing only a few trainable parameters and their optimizer state. Transformer
continual-learning methods such as DyTox reduce task-specific parameter growth
\cite{douillard2022dytox}, and bottleneck adapters achieve the same end by
inserting small trainable modules into a frozen backbone
\cite{houlsby2019adapters}. On the systems side, TinyTL observes that trainable
\emph{parameter} count is the wrong memory target because activations dominate
\cite{cai2020tinytl}, and on-device training has been demonstrated within a few
hundred kilobytes \cite{lin2022ondevice} on runtimes such as TensorFlow Lite
Micro \cite{david2021tflitemicro}, whose workloads the MLPerf Tiny suite
standardizes \cite{banbury2021mlperftiny}. The problem has stayed active: recent
TinyML continual-learning work adapts model \emph{size} to the task and distils
the retained set rather than storing raw exemplars \cite{ruib2024tinyml}, and a
2025 survey of replay-based continual learning on edge hardware confirms that
the binding constraint in practice is how much history fits in memory rather than
which replay rule is used \cite{replay2025survey} --- which is the premise this
paper takes to its conclusion by making the byte budget the independent variable.
Parameter-efficient adaptation of ECG foundation models, including LoRA, is now
surveyed as a standard branch \cite{ecgfoundation2024survey}, but at a parameter
scale far above the microcontroller regime studied here. Our accounting is the persistent-state
complement to that observation: we bound what must be \emph{retained} between
updates rather than what peaks during one. None of this work studies the
ECG-specific replay-depth discontinuity under a kilobyte-scale payload budget.

\subsection{Replay selection}

Replay performance depends on which exemplars are retained. Herding,
gradient-based selection, and bilevel coreset optimization are representative
approaches \cite{rebuffi2017icarl,tiwari2022gcr,borsos2020coresets}. Selection
quality matters, but our evidence shows that a prior question must be answered
first: \emph{how much of the budget should be spent on replay at all}, versus on
representation plasticity? We therefore hold the selection policy fixed and vary
the allocation between replay volume and trainable depth.

\subsection{Class imbalance in minority-beat classification}

Supraventricular and fusion beats are rare, and the source model's failure on
them is what patient adaptation must repair. Re-weighting by effective sample
number \cite{cui2019classbalanced} and focal loss \cite{lin2017focal} are the
standard remedies; we use both to build the class-reweighted source
variants of RQ5, in order to test whether our adaptation gain is merely
compensating for a weak source rather than adding anything.

\subsection{What is new here, and what is not}

We want to be explicit about the boundary. Low-rank adaptation
\cite{hu2022lora}, replay \cite{rebuffi2017icarl}, latent replay
\cite{pellegrini2020latent}, patient-specific ECG classification
\cite{hu1997patient,kiranyaz2016realtime}, and the tiny ECG transformer
architecture itself \cite{busia2025tiny,vaswani2017attention} are all
established, and none of them is claimed as a contribution of this paper.

What is new is their \emph{fixed-byte co-design}: measuring that replay cost is
non-monotonic in depth for a tiny ECG transformer, formulating a single
persistent-state budget that replay volume and trainable depth must share,
generating budget-matched configurations from that constraint, and evaluating the
resulting allocation under controlled patient-sequential adaptation. To our
knowledge the allocation question --- given a fixed number of bytes, how much
should be spent on remembering the source versus on the freedom to change the
representation --- has not been posed for this class of model.
Table~\ref{tab:relwork} places the present work against the closest prior
studies.

\begin{table}[H]
\centering
\small
\caption{Positioning against representative prior work. ``Budget'' is an explicit
byte constraint on retained adaptation state, not a parameter count.}
\label{tab:relwork}
\setlength{\tabcolsep}{3.5pt}
\begin{tabular}{lccll@{}c@{}}
\toprule
Work & ECG & Pt.-spec. & Replay & Trainable scope & Byte budget \\
\midrule
Hu \cite{hu1997patient}                  & yes & yes & ---            & patient expert  & no \\
de Chazal \cite{dechazal2004automatic}   & yes & no  & ---            & full model      & no \\
Kiranyaz \cite{kiranyaz2016realtime}     & yes & yes & ---            & full 1-D CNN    & no \\
Hannun \cite{hannun2019cardiologist}     & yes & no  & ---            & full model      & no \\
Pellegrini \cite{pellegrini2020latent}   & no  & no  & latent         & partial net     & no \\
Ravaglia \cite{ravaglia2021tinyml}       & no  & no  & quant.\ latent & partial net     & yes \\
Lin \cite{lin2022ondevice}               & no  & no  & ---            & sparse layers   & peak \\
Hu \cite{hu2022lora}                     & no  & no  & ---            & low-rank        & no \\
Busia \cite{busia2025tiny}               & yes & no  & ---            & inference only  & infer.\ \\
\midrule
\textbf{This work} & yes & yes & raw $+$ post-pool & head / adapter / LoRA
 & \textbf{persist.} \\
\bottomrule
\end{tabular}
\end{table}
